% !TeX program = xelatex
\documentclass[UTF8,a4paper,11pt]{ctexart}

\usepackage[left=2cm,right=2cm,top=2cm,bottom=2cm]{geometry}
\usepackage{amsmath,amssymb,bm}
\usepackage{booktabs}
\usepackage{float}
\usepackage{graphicx}
\usepackage{xcolor}
\usepackage{listings}
\usepackage{fancyhdr}
\usepackage[hidelinks]{hyperref}

% ======================== 考前只改这里 ========================
\newcommand{\CourseName}{数值分析期末上机}
\newcommand{\ProblemTitle}{题目一：经典四阶 Runge--Kutta 方法}
\newcommand{\StudentID}{241098117}
\newcommand{\StudentName}{张城宗}
\newcommand{\PythonFile}{单题Python代码模板.py}
% ============================================================

\setlength{\parindent}{2em}
\setlength{\parskip}{0.25em}
\setlength{\headheight}{14pt}
\pagestyle{fancy}
\fancyhf{}
\fancyhead[L]{\CourseName\quad 学号：\StudentID\quad 姓名：\StudentName}
\fancyhead[R]{\thepage}

\lstdefinestyle{python}{
  language=Python,
  basicstyle=\ttfamily\footnotesize,
  keywordstyle=\color{blue!75!black},
  commentstyle=\color{green!45!black},
  stringstyle=\color{orange!70!black},
  numbers=left,
  numberstyle=\tiny\color{gray},
  frame=single,
  breaklines=true,
  showstringspaces=false
}

% 图片不存在时显示占位框，因此可以先写报告、后运行 Python。
\newcommand{\resultfigure}[2]{
  \begin{figure}[H]
    \centering
    \IfFileExists{#1}{
      \includegraphics[width=0.78\textwidth]{#1}
    }{
      \fbox{\parbox{0.72\textwidth}{\centering
        \vspace{1.2em}运行 \texttt{\PythonFile} 后自动生成 \texttt{#1}\vspace{1.2em}}}
    }
    \caption{#2}
  \end{figure}
}

\begin{document}

\begin{center}
  {\LARGE\bfseries \CourseName 报告}

  \vspace{0.3em}
  学号：\StudentID\qquad 姓名：\StudentName
\end{center}

\section{\ProblemTitle}

\subsection{问题描述}

% 【必须替换】抄清楚方程、区间、初值、步长、要求。
求解初值问题
\[
  y'=f(x,y),\qquad x\in[x_0,X],\qquad y(x_0)=y_0,
\]
取步长 $h=\underline{\hspace{2cm}}$，使用
\underline{\hspace{4cm}} 方法计算数值解，并分析误差。

\subsection{方法与公式}

% 【必须替换】保留本题真正使用的方法，删除无关文字。
以经典四阶 Runge--Kutta 方法为例：
\[
  \begin{aligned}
    K_1&=f(x_n,y_n),\\
    K_2&=f\left(x_n+\frac h2,y_n+\frac h2K_1\right),\\
    K_3&=f\left(x_n+\frac h2,y_n+\frac h2K_2\right),\\
    K_4&=f(x_n+h,y_n+hK_3),\\
    y_{n+1}&=y_n+\frac h6(K_1+2K_2+2K_3+K_4).
  \end{aligned}
\]

若题目给出精确解 $y(x)$，使用最大绝对误差
\[
  E_\infty=\max_n|y_n-y(x_n)|
\]
评价数值结果。若需验证收敛阶，则使用
\[
  p\approx\frac{\log(E(h_1)/E(h_2))}{\log(h_1/h_2)}.
\]

\subsection{Python 实现说明}

程序首先建立等距网格，再按上述公式逐步推进；最后输出节点结果、
最大绝对误差和终点误差，并生成解曲线与误差曲线。

% 【必须替换】若有隐式迭代、起步方法或特殊停止条件，在这里说明。

\subsection{数值结果}

% Python 会生成 result_table.tex；未运行程序时显示提示。
\IfFileExists{result_table.tex}{
  \begin{table}[H]
  \centering
  \caption{各节点数值结果}
  \begin{tabular}{cccc}
    \toprule
    $x_n$ & 数值解 $y_n$ & 精确解 $y(x_n)$ & 绝对误差 \\
    \midrule
    0.000000 & 1.0000000000 & 1.0000000000 & 0.000e+00 \\
    0.100000 & 1.0003252083 & 1.0003251639 & 4.441e-08 \\
    0.200000 & 1.0025385940 & 1.0025384938 & 1.002e-07 \\
    0.300000 & 1.0083637234 & 1.0083635586 & 1.648e-07 \\
    0.400000 & 1.0193601439 & 1.0193599079 & 2.360e-07 \\
    0.500000 & 1.0369389926 & 1.0369386806 & 3.120e-07 \\
    0.600000 & 1.0623771190 & 1.0623767278 & 3.912e-07 \\
    0.700000 & 1.0968298648 & 1.0968293924 & 4.723e-07 \\
    0.800000 & 1.1413426261 & 1.1413420718 & 5.543e-07 \\
    0.900000 & 1.1968613168 & 1.1968606805 & 6.362e-07 \\
    1.000000 & 1.2642418350 & 1.2642411177 & 7.174e-07 \\
    \bottomrule
  \end{tabular}
\end{table}
}{
  \begin{center}
    \fbox{\parbox{0.72\textwidth}{\centering
      运行 \texttt{\PythonFile} 后自动插入数值结果表。}}
  \end{center}
}

\resultfigure{solution.pdf}{数值解与精确解比较}
\resultfigure{error.pdf}{各节点绝对误差}

\subsection{结果分析}

\begin{enumerate}
  \item \textbf{数值现象：}【必须替换】说明数值解的变化趋势，以及是否符合题意。
  \item \textbf{误差与精度：}【必须替换】报告最大误差、终点误差或估计收敛阶。
  \item \textbf{原因解释：}【必须替换】联系方法阶数、步长、稳定性或迭代容差解释结果。
\end{enumerate}

\subsection{结论}

【必须替换】用两三句话回答题目要求，并给出最重要的数值结论。

\clearpage
\section*{附录：Python 程序}
\addcontentsline{toc}{section}{附录：Python 程序}

\IfFileExists{\PythonFile}{
  \lstinputlisting[style=python,caption={本题 Python 程序}]{\PythonFile}
}{
  未找到 \texttt{\PythonFile}。
}

\end{document}
